%=====================================================================
\chapter{Overfitting, Bias--Variance and Regularisation}\label{ch:generalization}
%=====================================================================
\locator{2}{Part II --- Linear Models and Generalization}

\begin{outcomes}
By the end of this chapter you will be able to:
\begin{tight}
  \item \textbf{Distinguish} training error from true error, and
        \textbf{recognise} underfitting and overfitting from both.
  \item \textbf{Decompose} expected error into bias, variance and noise, and
        \textbf{compute} each from repeated predictions.
  \item \textbf{Diagnose} a model from its learning curves and \textbf{choose} the
        right remedy.
  \item \textbf{Apply} ridge and lasso regularisation and \textbf{explain} how the
        penalty strength trades bias for variance.
  \item \textbf{Select} a model and its hyperparameters with $k$-fold
        cross-validation, touching the test set only once.
\end{tight}
\end{outcomes}

\begin{prereq}
\chref{workflow} (training, validation and test sets; pipelines),
\chref{linreg} (polynomial features, MSE) and \chref{logreg} (the growing weights
on separable data).
\end{prereq}

\begin{keyterms}
generalisation $\bullet$ training error $\bullet$ true error $\bullet$
generalisation gap $\bullet$ underfitting $\bullet$ overfitting $\bullet$ model
complexity $\bullet$ bias $\bullet$ variance $\bullet$ irreducible error (noise)
$\bullet$ learning curve $\bullet$ regularisation $\bullet$ ridge (L2) $\bullet$
lasso (L1) $\bullet$ hyperparameter $\bullet$ $k$-fold cross-validation $\bullet$
grid search $\bullet$ nested cross-validation
\end{keyterms}

\section{The Goal Is Generalisation}

A model's \term{training error} is its error on the examples it learned from. Its
\term{true error} is its expected error on new examples drawn from the same
source --- the only error anybody cares about. The difference between the two is
the \term{generalisation gap}. Training error is easy to measure and always
optimistic; true error cannot be measured at all, only estimated on data the
model has never seen. Everything in this chapter follows from taking that
distinction seriously.

\section{Underfitting and Overfitting}

\dsref{polyfits} fits polynomials of three degrees to the same ten noisy
measurements of a smooth rise and fall: a team's weekly ticket volume over one
release cycle, which climbs after each release and falls as the fixes land ($x$ is
the position in the cycle, $y$ the volume on a scaled axis). Degree 1 is too rigid
to follow the curve at all. Degree 9 has ten
coefficients for ten points, so it passes through every one of them --- training
error exactly zero --- and swings wildly between them. Degree 3 follows the shape
and ignores the noise.

\dsfig{%
\begin{tikzpicture}
\begin{groupplot}[
  group style={group size=3 by 1, horizontal sep=0.9cm},
  width=5.6cm, height=4.6cm, axis lines=left,
  tick label style={font=\tiny}, label style={font=\tiny},
  title style={font=\scriptsize\bfseries},
  xmin=0, xmax=1, ymin=-2, ymax=2, xtick={0,0.5,1}, ytick={-2,-1,0,1,2}]
\nextgroupplot[title={Degree 1: underfits}, ylabel={$y$}]
\addplot[gray!55, dashed, domain=0:1, samples=60] {sin(deg(2*pi*x))};
\addplot[accent, line width=1.2pt] coordinates {(0.020,0.761) (0.030,0.746) (0.040,0.731) (0.050,0.716) (0.060,0.701) (0.070,0.686) (0.080,0.671) (0.090,0.656) (0.100,0.641) (0.110,0.627) (0.120,0.612) (0.130,0.597) (0.140,0.582) (0.150,0.567) (0.160,0.552) (0.170,0.537) (0.180,0.522) (0.190,0.507) (0.200,0.492) (0.210,0.477) (0.220,0.462) (0.230,0.447) (0.240,0.432) (0.250,0.417) (0.260,0.402) (0.270,0.387) (0.280,0.372) (0.290,0.357) (0.300,0.342) (0.310,0.327) (0.320,0.312) (0.330,0.297) (0.340,0.282) (0.350,0.267) (0.360,0.252) (0.370,0.237) (0.380,0.222) (0.390,0.207) (0.400,0.192) (0.410,0.177) (0.420,0.162) (0.430,0.147) (0.440,0.133) (0.450,0.118) (0.460,0.103) (0.470,0.088) (0.480,0.073) (0.490,0.058) (0.500,0.043) (0.510,0.028) (0.520,0.013) (0.530,-0.002) (0.540,-0.017) (0.550,-0.032) (0.560,-0.047) (0.570,-0.062) (0.580,-0.077) (0.590,-0.092) (0.600,-0.107) (0.610,-0.122) (0.620,-0.137) (0.630,-0.152) (0.640,-0.167) (0.650,-0.182) (0.660,-0.197) (0.670,-0.212) (0.680,-0.227) (0.690,-0.242) (0.700,-0.257) (0.710,-0.272) (0.720,-0.287) (0.730,-0.302) (0.740,-0.317) (0.750,-0.332) (0.760,-0.347) (0.770,-0.361) (0.780,-0.376) (0.790,-0.391) (0.800,-0.406) (0.810,-0.421) (0.820,-0.436) (0.830,-0.451) (0.840,-0.466) (0.850,-0.481) (0.860,-0.496) (0.870,-0.511) (0.880,-0.526) (0.890,-0.541) (0.900,-0.556) (0.910,-0.571) (0.920,-0.586) (0.930,-0.601) (0.940,-0.616) (0.950,-0.631) (0.960,-0.646) (0.970,-0.661) (0.980,-0.676) (0.990,-0.691)};
\addplot[only marks, mark=*, mark size=1.6pt, primary] coordinates {(0.034,0.313) (0.141,0.623) (0.262,0.482) (0.365,1.276) (0.433,0.409) (0.531,-0.078) (0.651,-0.974) (0.772,-1.197) (0.870,-0.567) (0.976,0.089)};
\nextgroupplot[title={Degree 3: about right}]
\addplot[gray!55, dashed, domain=0:1, samples=60] {sin(deg(2*pi*x))};
\addplot[greenhl!70!black, line width=1.2pt] coordinates {(0.020,-0.024) (0.030,0.083) (0.040,0.184) (0.050,0.278) (0.060,0.366) (0.070,0.447) (0.080,0.522) (0.090,0.591) (0.100,0.654) (0.110,0.711) (0.120,0.763) (0.130,0.809) (0.140,0.849) (0.150,0.885) (0.160,0.915) (0.170,0.941) (0.180,0.961) (0.190,0.977) (0.200,0.989) (0.210,0.996) (0.220,0.999) (0.230,0.999) (0.240,0.994) (0.250,0.985) (0.260,0.973) (0.270,0.958) (0.280,0.939) (0.290,0.917) (0.300,0.892) (0.310,0.864) (0.320,0.833) (0.330,0.800) (0.340,0.764) (0.350,0.726) (0.360,0.686) (0.370,0.644) (0.380,0.601) (0.390,0.555) (0.400,0.508) (0.410,0.459) (0.420,0.410) (0.430,0.359) (0.440,0.307) (0.450,0.254) (0.460,0.201) (0.470,0.147) (0.480,0.093) (0.490,0.039) (0.500,-0.016) (0.510,-0.070) (0.520,-0.124) (0.530,-0.178) (0.540,-0.232) (0.550,-0.284) (0.560,-0.336) (0.570,-0.387) (0.580,-0.437) (0.590,-0.486) (0.600,-0.533) (0.610,-0.579) (0.620,-0.623) (0.630,-0.665) (0.640,-0.706) (0.650,-0.744) (0.660,-0.780) (0.670,-0.813) (0.680,-0.844) (0.690,-0.873) (0.700,-0.898) (0.710,-0.920) (0.720,-0.940) (0.730,-0.956) (0.740,-0.968) (0.750,-0.978) (0.760,-0.983) (0.770,-0.984) (0.780,-0.982) (0.790,-0.975) (0.800,-0.964) (0.810,-0.949) (0.820,-0.928) (0.830,-0.904) (0.840,-0.874) (0.850,-0.839) (0.860,-0.799) (0.870,-0.754) (0.880,-0.703) (0.890,-0.647) (0.900,-0.585) (0.910,-0.517) (0.920,-0.443) (0.930,-0.362) (0.940,-0.276) (0.950,-0.183) (0.960,-0.083) (0.970,0.024) (0.980,0.137) (0.990,0.258)};
\addplot[only marks, mark=*, mark size=1.6pt, primary] coordinates {(0.034,0.313) (0.141,0.623) (0.262,0.482) (0.365,1.276) (0.433,0.409) (0.531,-0.078) (0.651,-0.974) (0.772,-1.197) (0.870,-0.567) (0.976,0.089)};
\nextgroupplot[title={Degree 9: overfits}]
\addplot[gray!55, dashed, domain=0:1, samples=60] {sin(deg(2*pi*x))};
\addplot[accent, line width=1.2pt] coordinates {(0.020,-3.000) (0.030,-1.382) (0.040,2.785) (0.050,3.000) (0.060,3.000) (0.070,3.000) (0.080,3.000) (0.090,3.000) (0.100,3.000) (0.110,3.000) (0.120,2.433) (0.130,1.484) (0.140,0.664) (0.150,-0.001) (0.160,-0.503) (0.170,-0.844) (0.180,-1.036) (0.190,-1.095) (0.200,-1.041) (0.210,-0.897) (0.220,-0.685) (0.230,-0.427) (0.240,-0.144) (0.250,0.148) (0.260,0.432) (0.270,0.695) (0.280,0.928) (0.290,1.123) (0.300,1.276) (0.310,1.384) (0.320,1.448) (0.330,1.470) (0.340,1.453) (0.350,1.401) (0.360,1.320) (0.370,1.217) (0.380,1.096) (0.390,0.965) (0.400,0.829) (0.410,0.693) (0.420,0.562) (0.430,0.441) (0.440,0.331) (0.450,0.236) (0.460,0.155) (0.470,0.089) (0.480,0.038) (0.490,-0.000) (0.500,-0.028) (0.510,-0.048) (0.520,-0.063) (0.530,-0.076) (0.540,-0.092) (0.550,-0.112) (0.560,-0.141) (0.570,-0.180) (0.580,-0.233) (0.590,-0.298) (0.600,-0.379) (0.610,-0.473) (0.620,-0.581) (0.630,-0.699) (0.640,-0.825) (0.650,-0.955) (0.660,-1.085) (0.670,-1.209) (0.680,-1.324) (0.690,-1.424) (0.700,-1.503) (0.710,-1.557) (0.720,-1.582) (0.730,-1.575) (0.740,-1.535) (0.750,-1.462) (0.760,-1.357) (0.770,-1.224) (0.780,-1.070) (0.790,-0.901) (0.800,-0.729) (0.810,-0.566) (0.820,-0.425) (0.830,-0.321) (0.840,-0.270) (0.850,-0.285) (0.860,-0.380) (0.870,-0.563) (0.880,-0.837) (0.890,-1.198) (0.900,-1.629) (0.910,-2.102) (0.920,-2.568) (0.930,-2.955) (0.940,-3.000) (0.950,-3.000) (0.960,-2.479) (0.970,-1.184) (0.980,1.100) (0.990,3.000)};
\addplot[only marks, mark=*, mark size=1.6pt, primary] coordinates {(0.034,0.313) (0.141,0.623) (0.262,0.482) (0.365,1.276) (0.433,0.409) (0.531,-0.078) (0.651,-0.974) (0.772,-1.197) (0.870,-0.567) (0.976,0.089)};
\end{groupplot}
\end{tikzpicture}}%
{Three polynomials fitted to the same ten points (blue); the dashed curve is the true
function. Training error falls from left to right and reaches zero at degree 9, which
is nonetheless the worst of the three on new data.}{polyfits}

\begin{definitionbox}[Underfitting and Overfitting]
A model \term{underfits} when it is too simple to capture the pattern: both
training and test error are high. A model \term{overfits} when it captures the
noise in its particular training sample as well as the pattern: training error is
low, test error is high, and the gap between them is large.
\end{definitionbox}

\dsfig{%
\begin{tikzpicture}
\begin{axis}[width=9.6cm, height=5.4cm, axis lines=left,
  xlabel={polynomial degree (model complexity)}, ylabel={RMSE},
  tick label style={font=\tiny}, label style={font=\scriptsize},
  xmin=0, xmax=9.3, ymin=0, ymax=1.6, xtick={0,...,9},
  legend style={font=\tiny, draw=gray!40, at={(0.03,0.97)}, anchor=north west}]
\fill[accent!6] (axis cs:0,0) rectangle (axis cs:2.5,1.6);
\fill[accent!6] (axis cs:7.5,0) rectangle (axis cs:9.3,1.6);
\addplot[primary, line width=1.2pt, mark=*, mark size=1.4pt] coordinates
  {(0,0.722)(1,0.568)(2,0.567)(3,0.290)(4,0.263)(5,0.206)(6,0.186)(7,0.169)(8,0.128)(9,0.000)};
\addlegendentry{training error}
\addplot[accent, line width=1.2pt, mark=square*, mark size=1.4pt] coordinates
  {(0,0.777)(1,0.509)(2,0.508)(3,0.248)(4,0.271)(5,0.298)(6,0.326)(7,0.329)(8,0.502)(9,1.500)};
\addlegendentry{test error}
\node[font=\tiny\bfseries, accent!80!black] at (axis cs:1.25,1.45) {underfitting};
\node[font=\tiny\bfseries, accent!80!black] at (axis cs:8.4,1.15) {overfitting};
\draw[greenhl!70!black, dashed, line width=0.9pt] (axis cs:3,0) -- (axis cs:3,1.0);
\node[font=\tiny, greenhl!45!black, anchor=south] at (axis cs:3,1.0) {best: degree 3};
\node[font=\tiny, accent, anchor=east] at (axis cs:8.85,1.5) {1.69 $\rightarrow$};
\end{axis}
\end{tikzpicture}}%
{Training and test error for the polynomials of \dsref{polyfits}, degree 0 to 9, with the
test error measured on 2{,}000 fresh samples. Training error only ever falls; test error
falls, bottoms out, and rises. The shape of that second curve is the whole of this
chapter.}{complexity-curve}

\section{The Bias--Variance Decomposition}

Why does test error rise again? Imagine repeating the whole experiment many
times: draw a fresh training set, fit the model, predict at one fixed input
$x_0$. The predictions scatter. Their average may miss the truth, and they may be
spread widely around their own average. Those are two different failings.

\begin{definitionbox}[Bias, Variance and Noise]
For a fixed input $x_0$ with true value $f(x_0)$, let $\hat{f}(x_0)$ be the
prediction of a model trained on a random training set. Then the expected squared
error on a new observation $y = f(x_0) + \varepsilon$ is
\[
\mathbb{E}\big[(y - \hat{f}(x_0))^2\big]
= \underbrace{\big(\mathbb{E}[\hat{f}(x_0)] - f(x_0)\big)^2}_{\text{bias}^2}
+ \underbrace{\mathbb{E}\big[(\hat{f}(x_0) - \mathbb{E}[\hat{f}(x_0)])^2\big]}_{\text{variance}}
+ \underbrace{\sigma^2_{\varepsilon}}_{\text{noise}}
\]
\term{Bias} is how far the average prediction is from the truth --- the error of
the model's assumptions. \term{Variance} is how much the prediction changes from
one training set to another --- its sensitivity to the particular sample.
\term{Noise} is the \term{irreducible error} in the labels themselves, which no
model can remove.
\end{definitionbox}

\dsfig{%
\begin{tikzpicture}[font=\scriptsize]
\foreach \cx/\cy/\ttl/\sx/\sy/\spread in {
    0/0/{Low bias, low variance}/0/0/0.18,
    4.2/0/{Low bias, high variance}/0/0/0.62,
    0/-4.0/{High bias, low variance}/0.7/0.45/0.18,
    4.2/-4.0/{High bias, high variance}/0.7/0.45/0.62}{
  \begin{scope}[shift={(\cx,\cy)}]
    \fill[accent!10] (0,0) circle (1.5);
    \fill[white] (0,0) circle (1.1);
    \fill[accent!18] (0,0) circle (0.7);
    \fill[white] (0,0) circle (0.35);
    \fill[accent!45] (0,0) circle (0.12);
    \draw[gray!50] (0,0) circle (1.5);
    \foreach \a/\r in {10/0.9,75/0.6,140/1.0,200/0.4,260/0.85,320/0.7,30/0.3,180/0.75}{
      \fill[primary] ({\sx+\spread*\r*cos(\a)},{\sy+\spread*\r*sin(\a)}) circle (1.6pt);}
    \node[font=\scriptsize\bfseries, text=primary, anchor=north] at (0,-1.6) {\ttl};
  \end{scope}}
\node[anchor=north west, align=left, text width=5.6cm, font=\scriptsize, text=gray!55!black]
     at (6.4,1.5)
     {\textbf{The centre is the truth $f(x_0)$.} Each dot is the prediction of a model
      trained on a different random training set.\\[4pt]
      \textbf{Bias} is how far the \emph{cloud's centre} is from the bull's-eye:
      simple models such as a straight line through a curve miss in the same
      direction every time.\\[4pt]
      \textbf{Variance} is how \emph{spread out} the cloud is: flexible models such
      as the degree-9 polynomial land somewhere different every time.\\[4pt]
      Increasing complexity moves a model from the bottom-left target towards the
      top-right one.};
\end{tikzpicture}}%
{Bias and variance as target practice. The goal is the top-left target; in practice
lowering one tends to raise the other.}{bias-variance}

\begin{worked}[Bias and Variance From Four Training Sets]
The true value at $x_0$ is $f(x_0) = 5$, and the label noise has variance
$\sigma^2_\varepsilon = 1$. Two models are each trained on the same four random
training sets and predict at $x_0$:
\begin{tight}
  \item \textbf{Simple model:} $3.8,\ 4.2,\ 4.0,\ 4.0$. Mean $4.0$.\\
        $\text{bias}^2 = (4.0 - 5)^2 = 1.00$;\quad
        variance $= \tfrac{0.04 + 0.04 + 0 + 0}{4} = 0.02$.
  \item \textbf{Complex model:} $6.5,\ 3.4,\ 5.9,\ 4.2$. Mean $5.0$.\\
        $\text{bias}^2 = 0$;\quad
        variance $= \tfrac{2.25 + 2.56 + 0.81 + 0.64}{4} = 1.565$.
\end{tight}
\textbf{Expected squared error} $=$ bias$^2$ $+$ variance $+$ noise:
\begin{tight}
  \item simple: $1.00 + 0.02 + 1 = \mathbf{2.02}$
  \item complex: $0 + 1.565 + 1 = \mathbf{2.57}$
\end{tight}
\smallskip
\textbf{The biased model wins.} It is wrong on average, but consistently so; the
unbiased model is right on average and wrong every single time. And neither can
go below 1.0 --- that is the noise.
\end{worked}

\begin{conceptbox}[The Trade-Off]
Making a model more flexible lowers its bias and raises its variance; making it
more rigid does the reverse. Test error is their sum plus noise, so it is
U-shaped in complexity (\dsref{complexity-curve}). The practical question is never
``how do I eliminate bias?'' or ``how do I eliminate variance?'' but ``where is the
bottom of the U, for this much data?'' --- because more data lowers variance and
moves the bottom to the right.
\end{conceptbox}

\section{Noisy Data}

Noise is not only measurement error in the features. It is also \emph{label}
noise: a task recorded as late that was in fact cancelled, a bug report closed as
``not a bug'' by a tired reviewer. Noise sets the floor on achievable error,
and it is exactly what an overfitting model memorises. Two consequences follow.
First, a training error far below the noise level is a symptom, not an
achievement --- the model has fitted the noise. Second, the noisier the data, the
simpler the model it will support: every method in Part III has a setting (tree
depth and pruning, $k$ in nearest neighbours, $C$ in an SVM) whose job is to stop
it chasing noise.

\section{Diagnosing With Learning Curves}

A \term{learning curve} plots training and validation error against the number of
training examples. Its shape tells you which of bias and variance is the problem,
and therefore what to do.

\dsfig{%
\begin{tikzpicture}
\begin{groupplot}[
  group style={group size=2 by 1, horizontal sep=1.3cm},
  width=6.8cm, height=4.6cm, axis lines=left,
  tick label style={font=\tiny}, label style={font=\tiny},
  title style={font=\scriptsize\bfseries},
  xlabel={training set size}, ylabel={error},
  xmin=0, xmax=100, ymin=0, ymax=1, xtick=\empty, ytick=\empty]
\nextgroupplot[title={High bias}]
\addplot[primary, line width=1.2pt, domain=3:100, samples=60] {0.62-0.35*exp(-x/12)};
\addplot[accent, line width=1.2pt, domain=3:100, samples=60] {0.66+0.3*exp(-x/12)};
\addplot[gray!55, dashed] coordinates {(0,0.2)(100,0.2)};
\node[font=\tiny, primary, anchor=north east] at (axis cs:98,0.6) {training};
\node[font=\tiny, accent, anchor=south east] at (axis cs:98,0.68) {validation};
\node[font=\tiny, gray!55!black, anchor=north east] at (axis cs:98,0.19) {acceptable error};
\node[font=\tiny, text=gray!55!black, anchor=north west, align=left] at (axis cs:38,0.47)
     {both high and close:\\more data will not help};
\nextgroupplot[title={High variance}]
\addplot[primary, line width=1.2pt, domain=3:100, samples=60] {0.02+0.12*(1-exp(-x/30))};
\addplot[accent, line width=1.2pt, domain=3:100, samples=60] {0.26+0.62*exp(-x/28)};
\addplot[gray!55, dashed] coordinates {(0,0.2)(100,0.2)};
\node[font=\tiny, primary, anchor=north east] at (axis cs:98,0.13) {training};
\node[font=\tiny, accent, anchor=south east] at (axis cs:98,0.28) {validation};
\node[font=\tiny, text=gray!55!black, anchor=north west, align=left] at (axis cs:35,0.75)
     {large gap, still closing:\\more data will help};
\end{groupplot}
\end{tikzpicture}}%
{Learning curves. With high bias the curves meet early at an unacceptable level; with
high variance there is a wide gap that more data would narrow.}{learning-curves}

\rowcolors{2}{white}{lightbg}
\begin{longtable}{@{}L{3.0cm}L{4.4cm}L{7.6cm}@{}}
\toprule
\rowcolor{primary!15}
\textbf{Diagnosis} & \textbf{Symptom} & \textbf{What helps} \\
\midrule
\endhead
High bias (underfitting) & Training error high; validation close to it & A more
flexible model; more or better features; less regularisation. More data does
\emph{not} help \\
High variance (overfitting) & Training error low; validation much higher & More
data; fewer features; a simpler model; more regularisation; ensembles that
average (\chref{ensembles}) \\
Both fine but still poor & Both curves meet at the noise level & Better data:
cleaner labels, more informative features. The algorithm is not the problem \\
\bottomrule
\end{longtable}

\section{Regularisation}

Choosing a lower polynomial degree is one way to control complexity, but a crude
one: the degree jumps in whole steps. \term{Regularisation} controls it smoothly,
by keeping the high-degree terms but charging the model for large weights.

\begin{definitionbox}[Ridge and Lasso]
Add a penalty on the size of the weights to the training loss:
\[
\text{Ridge (L2):}\quad \text{MSE} + \lambda\sum_{j} w_j^2
\qquad\qquad
\text{Lasso (L1):}\quad \text{MSE} + \lambda\sum_{j} |w_j|
\]
The \term{hyperparameter} $\lambda \ge 0$ sets the strength: $\lambda = 0$ is
ordinary least squares; a very large $\lambda$ forces every weight to 0 (the bias
$b$ is not penalised). Ridge shrinks all weights smoothly; lasso drives some of
them to exactly zero, so it also \emph{selects features}. Features must be
standardised first, or the penalty falls unequally on them.
\end{definitionbox}

\begin{worked}[Ridge Shrinkage in One Dimension]
Centred data $x = -2, -1, 0, 1, 2$ and $y = -3, -1, 0, 2, 2$, fitted by $\hat{y} = wx$.
Minimising $\sum(y_i - wx_i)^2 + \lambda w^2$ gives
\[
w = \frac{\sum x_i y_i}{\sum x_i^2 + \lambda}
\qquad\text{with}\quad \textstyle\sum x_i y_i = 6 + 1 + 0 + 2 + 4 = 13,\quad \sum x_i^2 = 10
\]
\begin{tight}
  \item $\lambda = 0$: $w = 13/10 = \mathbf{1.30}$ (least squares).
  \item $\lambda = 2$: $w = 13/12 = \mathbf{1.08}$.
  \item $\lambda = 10$: $w = 13/20 = \mathbf{0.65}$.
\end{tight}
The penalty adds $\lambda$ to the denominator, so it pulls the weight towards zero
--- a little when the data is plentiful ($\sum x^2$ large), a lot when it is not.
That is exactly the behaviour wanted: shrink hardest where the data says least.
\end{worked}

\dsfig{%
\begin{tikzpicture}[font=\scriptsize]
\foreach \ox/\ttl in {0/{Ridge (L2): a circle}, 6.6/{Lasso (L1): a diamond}}{
  \begin{scope}[shift={(\ox,0)}]
  \draw[-{Stealth[length=2mm]}, gray!60] (-2.2,0) -- (2.8,0) node[right, font=\tiny] {$w_1$};
  \draw[-{Stealth[length=2mm]}, gray!60] (0,-1.6) -- (0,3.3) node[above, font=\tiny] {$w_2$};
  \node[font=\scriptsize\bfseries, text=primary] at (0.3,-2.0) {\ttl};
  \end{scope}}
% ridge: least squares at (1.45, 1.2); tangent level 1.363, touching at (0.669, 0.743)
\fill[greenhl!18, draw=greenhl!70!black, line width=1pt] (0,0) circle (1.0);
\foreach \lv in {0.35,0.8,1.363}{
  \draw[secondary!70, line width=0.8pt, rotate around={-35:(1.45,1.2)}]
       (1.45,1.2) ellipse ({1.2*\lv} and {0.62*\lv});}
\fill[secondary] (1.45,1.2) circle (1.6pt);
\node[font=\tiny, text=secondary!70!black, anchor=south west] at (1.5,1.25) {least squares};
\fill[greenhl!40!black] (0.669,0.743) circle (1.9pt);
\node[font=\tiny, text=greenhl!40!black, anchor=north west, align=left] at (0.9,-0.1)
     {first contact:\\both weights shrunk,\\neither is 0};
\draw[greenhl!40!black, line width=0.4pt] (0.95,-0.12) -- (0.70,0.70);
% lasso: least squares at (0.55, 1.95); tangent level 1.632, touching the corner (0, 1.1)
\fill[greenhl!18, draw=greenhl!70!black, line width=1pt]
     (6.6+1.1,0) -- (6.6,1.1) -- (6.6-1.1,0) -- (6.6,-1.1) -- cycle;
\foreach \lv in {0.4,0.95,1.632}{
  \draw[secondary!70, line width=0.8pt, rotate around={-35:(7.15,1.95)}]
       (7.15,1.95) ellipse ({1.2*\lv} and {0.62*\lv});}
\fill[secondary] (7.15,1.95) circle (1.6pt);
\node[font=\tiny, text=secondary!70!black, anchor=south west] at (7.2,2.0) {least squares};
\fill[accent] (6.6,1.1) circle (2pt);
\node[font=\tiny, text=accent, anchor=north east, align=right] at (5.4,0.95)
     {first contact at a corner:\\$w_1 = 0$ exactly};
\draw[accent, line width=0.4pt] (5.42,0.92) -- (6.52,1.08);
\end{tikzpicture}}%
{Why lasso sets weights to zero. The regularised solution is where the smallest
error ellipse (blue) first touches the region allowed by the penalty (green). A circle
is touched at a generic point; a diamond is usually touched at a corner, where one
weight is exactly zero.}{ridge-lasso}

\begin{alertbox}[Regularisation Is Everywhere, Under Other Names]
Every flexible model in this course has a knob that trades variance for bias.
Maximum tree depth and pruning (\chref{trees}); $k$ in $k$-nearest neighbours
(\chref{knn}); $C$ in SVMs (\chref{svm}); weight decay, early stopping and dropout
in neural networks (\chref{ann}); $C = 1/\lambda$ in scikit-learn's
\texttt{LogisticRegression} (\chref{logreg}). Once you see them as the same idea,
tuning them is the same procedure --- the one in the next section.
\end{alertbox}

\section{Model Selection With Cross-Validation}

A \term{hyperparameter} is a setting chosen before training --- degree, $\lambda$,
$k$, tree depth --- as opposed to a \emph{parameter} such as a weight, which
training sets. Hyperparameters are chosen by validation error. With little data a
single validation split is too noisy to trust, so it is replaced by
cross-validation.

\begin{definitionbox}[$k$-Fold Cross-Validation]
Split the training data into $k$ equal folds (typically 5 or 10). For each fold in
turn, train on the other $k - 1$ folds and measure error on the held-out fold.
The cross-validation score is the mean of the $k$ errors; report their standard
deviation too. Every example is used for validation exactly once.
\end{definitionbox}

\dsfig{%
\begin{tikzpicture}[font=\scriptsize]
\foreach \r in {1,...,5}{
  \pgfmathsetmacro{\yy}{-(\r-1)*0.62}
  \node[anchor=east, font=\tiny, text=gray!55!black] at (-0.2,\yy) {round \r};
  \foreach \c in {1,...,5}{
    \pgfmathsetmacro{\xx}{(\c-1)*1.9+0.95}
    \ifnum\c=\r
      \node[rounded corners=2pt, fill=goldhl!30, draw=goldhl, minimum width=1.8cm,
            minimum height=0.48cm, font=\tiny\bfseries, text=goldhl!45!black] at (\xx,\yy) {validate};
    \else
      \node[rounded corners=2pt, fill=primary!12, draw=primary!50, minimum width=1.8cm,
            minimum height=0.48cm, font=\tiny, text=primary!55!black] at (\xx,\yy) {train};
    \fi}
  \node[anchor=west, font=\tiny, text=gray!55!black] at (9.6,\yy) {error$_\r$};}
\node[rounded corners=2pt, fill=gray!12, draw=gray!50, minimum width=1.8cm,
      minimum height=2.95cm, font=\tiny\bfseries, text=gray!45!black, align=center]
      at (12.4,-1.24) {TEST\\[2pt]set aside\\before\\any of this};
\node[anchor=north west, align=left, text width=11cm, font=\tiny, text=gray!55!black]
     at (-1.0,-2.95)
     {CV score $=$ mean of the five errors. Repeat for every candidate setting, pick the best,
      refit on all the training data, then evaluate on the test set once.};
\end{tikzpicture}}%
{Five-fold cross-validation on the training data. The test set is not part of the
rotation.}{kfold}

\Needspace{14\baselineskip}
\begin{worked}[Choosing $\lambda$ by Cross-Validation]
Five-fold cross-validation of a degree-12 polynomial with ridge penalties,
from this chapter's lab:

\smallskip
\begin{center}
\begin{tabular}{@{}lrl@{}}
\toprule
$\lambda$ & \textbf{CV RMSE} & \\
\midrule
$10^{-6}$ & 0.388 & \\
$10^{-3}$ & 0.326 & $\leftarrow$ \textbf{lowest} \\
$10^{-1}$ & 0.330 & \\
$10$ & 0.441 & \\
$1000$ & 0.706 & \\
\bottomrule
\end{tabular}
\end{center}

\smallskip
The U-shape again: too little penalty overfits, too much underfits. Choose
$\lambda = 10^{-3}$, refit on the whole training set, and only then measure the
test set: RMSE 0.299. Unregularised, the same degree-12 polynomial had a CV RMSE of
2.14 --- regularisation, not a lower degree, rescued it.

\textbf{Notice} that $10^{-3}$ and $10^{-1}$ are within 0.004 of each other, far less
than the fold-to-fold variation. When candidates are this close, prefer the
simpler model --- here the larger penalty.
\end{worked}

\begin{alertbox}[Tuning on the Test Set Is Still Cheating When a Library Does It]
If you try thirty settings and report the best \emph{test} score, the test set has
become a validation set and the score is optimistic. Tune with cross-validation on
the training data (\texttt{GridSearchCV} does this), then use the test set once.
When even the tuning procedure must be evaluated honestly --- for a published
comparison --- wrap the whole search in an outer cross-validation loop:
\term{nested cross-validation}.
\end{alertbox}

\begin{pitfall}
\begin{tight}
  \item \textbf{Celebrating zero training error.} With noisy labels it proves
        overfitting.
  \item \textbf{Collecting more data to cure high bias.} If training error is
        already high, more of the same data will not lower it.
  \item \textbf{Regularising unscaled features.} The penalty then depends on the
        units each feature happens to be measured in.
  \item \textbf{Reporting the best of many test scores.} Select on
        cross-validation; test once.
  \item \textbf{Choosing between near-identical CV scores.} Differences smaller
        than the fold-to-fold standard deviation are noise.
\end{tight}
\end{pitfall}

%---------------------------------------------------------------------
\begin{lab}[Find the Bottom of the U]
\begin{lstlisting}[language=Python]
import numpy as np
from sklearn.pipeline import make_pipeline
from sklearn.preprocessing import PolynomialFeatures, StandardScaler
from sklearn.linear_model import LinearRegression, Ridge
from sklearn.model_selection import (train_test_split, KFold,
                                     cross_val_score, GridSearchCV)

# --- a release cycle: ticket volume rises and falls, plus noise ---------
rng = np.random.default_rng(0)
X = rng.uniform(0, 1, (60, 1))
y = np.sin(2 * np.pi * X[:, 0]) + rng.normal(0, 0.25, 60)
X_tr, X_te, y_tr, y_te = train_test_split(X, y, test_size=0.25,
                                          random_state=0)
cv = KFold(5, shuffle=True, random_state=0)
rmse = "neg_root_mean_squared_error"

# --- 1. complexity: training error against cross-validated error ---------
for deg in (1, 3, 6, 12):
    m = make_pipeline(PolynomialFeatures(deg), StandardScaler(),
                      LinearRegression())
    m.fit(X_tr, y_tr)
    train = np.sqrt(np.mean((m.predict(X_tr) - y_tr) ** 2))
    cvs = -cross_val_score(m, X_tr, y_tr, cv=cv, scoring=rmse)
    print(f"degree {deg:2d}: train RMSE {train:.3f}   "
          f"CV RMSE {cvs.mean():.3f} +/- {cvs.std():.3f}")

# --- 2. keep degree 12, but regularise; choose alpha by CV ---------------
pipe = make_pipeline(PolynomialFeatures(12), StandardScaler(), Ridge())
grid = GridSearchCV(pipe, {"ridge__alpha": [1e-6, 1e-3, 1e-1, 10, 1000]},
                    cv=cv, scoring=rmse).fit(X_tr, y_tr)
for a, s in zip(grid.cv_results_["param_ridge__alpha"],
                grid.cv_results_["mean_test_score"]):
    print(f"alpha {a:>8g}: CV RMSE {-s:.3f}")
print("chosen:", grid.best_params_)

# --- 3. only now, once: the test set -------------------------------------
pred = grid.predict(X_te)
print("test RMSE:", round(float(np.sqrt(np.mean((pred - y_te) ** 2))), 3))
\end{lstlisting}
scikit-learn calls $\lambda$ \texttt{alpha}. Change the 60 samples to 600 and run
again: which degrees stop overfitting, and why?
\end{lab}

%---------------------------------------------------------------------
\begin{checkpoint}
\begin{tightnum}
  \item A model has training RMSE 0.05 and validation RMSE 0.90. Is the problem
        bias or variance? Name two remedies.
  \item Ten models trained on different samples predict $2, 2, 2, 2, 2, 2, 2, 2,
        2, 2$ at a point whose true value is 3. Give the bias$^2$ and the
        variance.
  \item Why does lasso, but not ridge, perform feature selection?
\end{tightnum}
\end{checkpoint}

%---------------------------------------------------------------------
\begin{chaptersummary}
\begin{tight}
  \item Training error is optimistic; generalisation is judged on unseen data.
  \item Underfitting: both errors high. Overfitting: training error low, test
        error high. Test error is U-shaped in model complexity.
  \item Expected error $=$ bias$^2$ $+$ variance $+$ noise. Flexibility trades bias
        for variance; noise is a floor no model can go below.
  \item Learning curves diagnose the problem: converged and high means bias; a
        wide, closing gap means variance.
  \item Ridge (L2) shrinks weights; lasso (L1) also sets some to zero. The
        penalty strength is a hyperparameter.
  \item Choose hyperparameters by $k$-fold cross-validation on the training
        data, refit, then test once.
\end{tight}
\end{chaptersummary}

\begin{reviewq}
  \item \textbf{(MCQ)} Adding training data is most useful when a model has:
        (a)~high bias (b)~high variance (c)~high noise (d)~low variance.
  \item \textbf{(MCQ)} As $\lambda$ in ridge regression increases: (a)~bias and
        variance both fall (b)~bias rises and variance falls (c)~bias falls and
        variance rises (d)~neither changes.
  \item \textbf{(MCQ)} Which of these is a hyperparameter? (a)~a regression
        weight (b)~the bias term $b$ (c)~the number of folds' mean error
        (d)~the degree of a polynomial.
  \item \textbf{(Short)} Explain why training error cannot be used to choose the
        degree of a polynomial.
  \item \textbf{(Short)} What is irreducible error, and why can no model remove
        it?
  \item \textbf{(Short)} Describe nested cross-validation and when it is needed.
  \item \textbf{(Applied)} Three models predict at a point with true value 10 and
        noise variance 4. Over five training sets they predict: A: $8, 8, 9, 8, 7$;
        B: $12, 7, 10, 13, 8$; C: $10, 10, 11, 9, 10$. Compute bias$^2$, variance and
        expected error for each, and rank them.
  \item \textbf{(Applied)} For the centred data $x = -1, 0, 1, 2$ (centre it first)
        and $y = 1, 2, 2, 5$ (centre it too), compute the ridge weight for
        $\lambda = 0, 1, 5$ and describe the pattern.
\end{reviewq}
